\documentclass[11pt]{article}

\usepackage[margin=1in]{geometry}
\usepackage{amsmath}
\usepackage{amssymb}
\usepackage{bm}
\usepackage{siunitx}
\usepackage{microtype}
\usepackage{hyperref}

\title{Relationship Between a Discrete Lattice Random Walk and Continuous Diffusion}
\author{}
\date{}

\begin{document}

\maketitle

\section{Overview}

The lattice Monte Carlo model represents molecular diffusion by allowing each particle to move between neighboring lattice sites with specified probabilities during a discrete timestep. Although the particle trajectories are restricted to discrete positions and times, the movement probabilities are chosen so that the model reproduces the continuous diffusion equation at sufficiently large spatial and temporal scales.

For a cubic lattice with lattice spacing \(a\), diffusion coefficient \(D\), and timestep \(\Delta t\), the probability of moving in each of the six Cartesian directions is

\begin{equation}
    p_{\mathrm{axis}}
    =
    \frac{D\Delta t}{a^2}.
\end{equation}

The probability of remaining at the current lattice site is

\begin{equation}
    p_{\mathrm{stay}}
    =
    1 - 6p_{\mathrm{axis}}
    =
    1-\frac{6D\Delta t}{a^2}.
\end{equation}

These probabilities arise directly from a finite-difference discretization of the continuous diffusion equation.

\section{Continuous Diffusion Equation}

Let

\begin{equation}
    c(x,y,z,t)
\end{equation}

denote either the concentration of diffusing particles or the probability density for the position of a single particle. In the absence of drift, reactions, and external forces, diffusion is governed by Fick's second law:

\begin{equation}
    \frac{\partial c}{\partial t}
    =
    D\nabla^2 c,
\end{equation}

where

\begin{equation}
    \nabla^2c
    =
    \frac{\partial^2c}{\partial x^2}
    +
    \frac{\partial^2c}{\partial y^2}
    +
    \frac{\partial^2c}{\partial z^2}.
\end{equation}

The diffusion coefficient \(D\) has units of area per unit time, such as
\(\si{\meter\squared\per\second}\).

\section{Finite-Difference Representation}

Consider a cubic lattice with spacing \(a\). Let

\begin{equation}
    c_{i,j,k}^{n}
\end{equation}

denote the probability or concentration at lattice location \((i,j,k)\) after timestep \(n\).

The second derivative in the \(x\)-direction can be approximated using a centered finite difference:

\begin{equation}
    \frac{\partial^2c}{\partial x^2}
    \approx
    \frac{
        c_{i+1,j,k}^{n}
        -2c_{i,j,k}^{n}
        +c_{i-1,j,k}^{n}
    }{a^2}.
\end{equation}

Applying the same approximation in the \(y\)- and \(z\)-directions gives

\begin{equation}
    \nabla^2c
    \approx
    \frac{1}{a^2}
    \left[
        \sum_{\text{six neighbors}}
        c_{\mathrm{neighbor}}^{n}
        -
        6c_{i,j,k}^{n}
    \right].
\end{equation}

The time derivative is approximated as

\begin{equation}
    \frac{\partial c}{\partial t}
    \approx
    \frac{
        c_{i,j,k}^{n+1}
        -
        c_{i,j,k}^{n}
    }{\Delta t}.
\end{equation}

Substitution into the diffusion equation gives

\begin{equation}
    \frac{
        c_{i,j,k}^{n+1}
        -
        c_{i,j,k}^{n}
    }{\Delta t}
    =
    \frac{D}{a^2}
    \left[
        \sum_{\text{six neighbors}}
        c_{\mathrm{neighbor}}^{n}
        -
        6c_{i,j,k}^{n}
    \right].
\end{equation}

Rearranging,

\begin{equation}
    c_{i,j,k}^{n+1}
    =
    \left(
        1-\frac{6D\Delta t}{a^2}
    \right)
    c_{i,j,k}^{n}
    +
    \frac{D\Delta t}{a^2}
    \sum_{\text{six neighbors}}
    c_{\mathrm{neighbor}}^{n}.
\end{equation}

Defining

\begin{equation}
    \lambda
    =
    \frac{D\Delta t}{a^2},
\end{equation}

the update becomes

\begin{equation}
    c_{i,j,k}^{n+1}
    =
    (1-6\lambda)c_{i,j,k}^{n}
    +
    \lambda
    \sum_{\text{six neighbors}}
    c_{\mathrm{neighbor}}^{n}.
\end{equation}

This expression has a direct probabilistic interpretation:

\begin{itemize}
    \item The particle moves to each neighboring lattice site with probability
    \(\lambda\).
    \item The particle remains at its current site with probability
    \(1-6\lambda\).
\end{itemize}

Therefore,

\begin{equation}
    p_{\mathrm{axis}}
    =
    \lambda
    =
    \frac{D\Delta t}{a^2}.
\end{equation}

The Monte Carlo transition probabilities are thus a stochastic implementation of the finite-difference diffusion equation.

\section{Derivation from Mean-Squared Displacement}

The same probability can be derived by requiring that each lattice step have the same displacement variance as continuous Brownian diffusion.

During one timestep, the displacement in the \(x\)-direction is

\begin{equation}
    \Delta X
    =
    \begin{cases}
        +a, & \text{with probability }p_{\mathrm{axis}},\\
        -a, & \text{with probability }p_{\mathrm{axis}},\\
        0,  & \text{otherwise}.
    \end{cases}
\end{equation}

Because the positive and negative directions are equally likely,

\begin{equation}
    \left\langle \Delta X \right\rangle
    =
    0.
\end{equation}

The mean-squared displacement in \(x\) during one timestep is

\begin{align}
    \left\langle \Delta X^2 \right\rangle
    &=
    p_{\mathrm{axis}}a^2
    +
    p_{\mathrm{axis}}a^2 \\
    &=
    2p_{\mathrm{axis}}a^2.
\end{align}

Continuous one-dimensional diffusion requires

\begin{equation}
    \left\langle \Delta X^2 \right\rangle
    =
    2D\Delta t.
\end{equation}

Equating the discrete and continuous expressions gives

\begin{equation}
    2p_{\mathrm{axis}}a^2
    =
    2D\Delta t,
\end{equation}

and therefore

\begin{equation}
    \boxed{
        p_{\mathrm{axis}}
        =
        \frac{D\Delta t}{a^2}
    }.
\end{equation}

The total three-dimensional mean-squared displacement during one timestep is

\begin{align}
    \left\langle \Delta R^2 \right\rangle
    &=
    \left\langle \Delta X^2 \right\rangle
    +
    \left\langle \Delta Y^2 \right\rangle
    +
    \left\langle \Delta Z^2 \right\rangle \\
    &=
    6D\Delta t.
\end{align}

After \(n\) independent timesteps,

\begin{equation}
    t=n\Delta t,
\end{equation}

and the displacement is the sum of the individual increments. Consequently,

\begin{align}
    \left\langle R^2(t) \right\rangle
    &=
    n\left(6D\Delta t\right) \\
    &=
    6Dt.
\end{align}

Thus, the lattice model reproduces the theoretical mean-squared displacement exactly in expectation.

\section{Emergence of the Gaussian Displacement Distribution}

A single lattice step is not Gaussian because only seven outcomes are possible: movement in one of six directions or no movement. After many timesteps, however, each coordinate displacement is the sum of many independent increments:

\begin{equation}
    X_n
    =
    \sum_{m=1}^{n}\Delta X_m.
\end{equation}

Each increment has zero mean and finite variance. The central limit theorem therefore implies that, after many steps, the coordinate displacement approaches a Gaussian distribution:

\begin{equation}
    X(t)
    \sim
    \mathcal{N}(0,2Dt).
\end{equation}

The same result holds for \(Y(t)\) and \(Z(t)\). The resulting three-dimensional displacement probability density is

\begin{equation}
    p(x,y,z,t)
    =
    \frac{1}{(4\pi Dt)^{3/2}}
    \exp\left[
        -\frac{x^2+y^2+z^2}{4Dt}
    \right].
\end{equation}

Equivalently, using

\begin{equation}
    r^2=x^2+y^2+z^2,
\end{equation}

the spatial probability density is

\begin{equation}
    p(\bm{r},t)
    =
    \frac{1}{(4\pi Dt)^{3/2}}
    \exp\left(
        -\frac{r^2}{4Dt}
    \right).
\end{equation}

The corresponding radial displacement probability density is obtained by multiplying by the volume of a spherical shell:

\begin{equation}
    p_R(r,t)
    =
    4\pi r^2
    \frac{1}{(4\pi Dt)^{3/2}}
    \exp\left(
        -\frac{r^2}{4Dt}
    \right),
    \qquad r\geq0.
\end{equation}

Therefore, the lattice displacement distribution becomes increasingly similar to the continuous Brownian propagator as the number of timesteps increases.

\section{Continuum Limit}

The formal relationship between the lattice model and continuous diffusion is obtained by considering the limit

\begin{equation}
    a\rightarrow0,
    \qquad
    \Delta t\rightarrow0,
\end{equation}

while keeping

\begin{equation}
    \frac{a^2}{\Delta t}
\end{equation}

scaled so that \(D\) remains constant.

A Taylor expansion of the concentrations at the neighboring lattice sites gives

\begin{equation}
    \sum_{\text{six neighbors}}
    c_{\mathrm{neighbor}}
    -
    6c
    =
    a^2\nabla^2c
    +
    \mathcal{O}(a^4).
\end{equation}

The lattice evolution equation can therefore be written as

\begin{equation}
    \frac{c^{n+1}-c^n}{\Delta t}
    =
    D\nabla^2c
    +
    \mathcal{O}(a^2)
    +
    \mathcal{O}(\Delta t).
\end{equation}

As \(a\) and \(\Delta t\) approach zero, the discretization errors vanish and the lattice equation converges to

\begin{equation}
    \frac{\partial c}{\partial t}
    =
    D\nabla^2c.
\end{equation}

This relationship is commonly described as the diffusive scaling limit or Brownian scaling limit of a random walk.

\section{Timestep Restriction}

All Monte Carlo transition probabilities must be nonnegative. In particular,

\begin{equation}
    p_{\mathrm{stay}}
    =
    1-\frac{6D\Delta t}{a^2}
    \geq0.
\end{equation}

Therefore, the timestep must satisfy

\begin{equation}
    \boxed{
        \Delta t
        \leq
        \frac{a^2}{6D}
    }.
\end{equation}

This is also the stability condition for the explicit finite-difference solution of the three-dimensional diffusion equation.

For example, if the simulation selects

\begin{equation}
    \Delta t
    =
    0.95\frac{a^2}{6D},
\end{equation}

then

\begin{equation}
    p_{\mathrm{axis}}
    =
    \frac{0.95}{6}
    \approx
    0.1583
\end{equation}

for each direction, and

\begin{equation}
    p_{\mathrm{stay}}
    =
    0.05.
\end{equation}

The total probability of moving during one timestep is \(0.95\), while the probability of remaining stationary is \(0.05\). The presence of stationary steps does not reduce the intended diffusion coefficient because the movement probabilities were calculated with the stationary probability included.

\section{Why the Numerical Agreement Is Close}

The close agreement between the lattice simulation and continuous diffusion theory follows from several properties of the model:

\begin{enumerate}
    \item The movement probability is selected to reproduce the Brownian displacement variance:
    \[
        p_{\mathrm{axis}}a^2
        =
        D\Delta t.
    \]

    \item The mean-squared displacement is therefore exactly
    \[
        \left\langle R^2(t)\right\rangle
        =
        6Dt
    \]
    in expectation.

    \item After many independent timesteps, the central limit theorem causes the coordinate displacement distributions to approach Gaussian distributions.

    \item A large particle ensemble reduces Monte Carlo sampling noise.

    \item In a free-diffusion validation protocol, boundaries, receptor reactions, ligand exchange, and external forces are absent and therefore cannot perturb the Brownian statistics.
\end{enumerate}

The lattice model does not merely approximate the slope of the theoretical mean-squared displacement empirically. The correct slope is encoded directly in the transition probabilities.

\section{Differences Between the Lattice and Continuous Descriptions}

Despite their close agreement at larger scales, the lattice and continuous models are not identical at all length and time scales.

\subsection{Short-Time Discreteness}

After one timestep, the radial displacement can only be

\begin{equation}
    r=0
    \qquad\text{or}\qquad
    r=a.
\end{equation}

Continuous diffusion permits every possible displacement. The discrete structure is therefore most apparent at very short times.

\subsection{Lattice Anisotropy}

A cubic lattice has preferred Cartesian directions. For example, a direct diagonal movement is not possible in one timestep. Approximate rotational symmetry emerges only after many steps.

\subsection{Finite Step Length}

The maximum displacement during one timestep is \(a\). In contrast, the mathematical diffusion propagator assigns a nonzero, although potentially extremely small, probability density at arbitrarily large distances for every \(t>0\).

\subsection{Absence of Inertial Dynamics}

Real particles may exhibit ballistic motion at sufficiently short times:

\begin{equation}
    \left\langle R^2(t)\right\rangle
    \propto
    t^2.
\end{equation}

The lattice model begins directly in the overdamped diffusive regime:

\begin{equation}
    \left\langle R^2(t)\right\rangle
    \propto
    t.
\end{equation}

This approximation is appropriate when the simulated timescale is much longer than the particle's inertial relaxation time.

\section{Conclusion}

The relationship between the lattice movement probabilities and continuous diffusion can be summarized as

\begin{equation}
    \boxed{
        p_{\mathrm{axis}}
        =
        \frac{
            \text{desired Brownian variance per coordinate and timestep}
        }{
            2\times\text{squared lattice spacing}
        }
        =
        \frac{D\Delta t}{a^2}
    }.
\end{equation}

This choice ensures that the discrete model has the same mean and variance as Brownian diffusion during each timestep. Repeated independent lattice movements then approach the Gaussian displacement distribution predicted by the continuous diffusion equation. The resulting lattice Monte Carlo model is therefore a consistent discrete representation of overdamped molecular diffusion, provided that the lattice spacing and timestep adequately resolve the physical scales of interest.

\end{document}